\section{Scale, Emergence y Alignment: de dónde vienen las capacidades}

Una vez entendida la attention, el curso amplía la perspectiva del mecanismo de una sola capa al entrenamiento a gran escala. Aquí hay tres líneas que no deben confundirse: cómo la escala modifica las capacidades observables, cómo la alignment modifica las preferencias de salida y cómo los modelos de producto adoptan distintas decisiones de compromiso en datos, arquitectura y receta de entrenamiento.

\subsection{Un Large Language Model es un sistema de entrenamiento ampliado}

Un Large language model (modelo de lenguaje grande, LLM) suele seguir teniendo como núcleo la next-token prediction, pero con una escala de parámetros, datos y cómputo considerablemente mayor. Su «large» no se refiere solo a los parámetros: también implica que crecen en conjunto el pipeline de datos, el entrenamiento en paralelo, la evaluación y la infraestructura de despliegue.

\lecturefigure{slide-028.jpg}{Escala, datos de entrenamiento y next-token objective de un LLM}{Diapositivas oficiales de CS25 V4, página 28}

\begin{importantbox}{Autoregressive objective}
Para una secuencia $x_{1:T}$, la pérdida de entrenamiento suele escribirse como
\[
\mathcal{L}_{\mathrm{LM}}(\theta)=-\sum_{t=1}^{T}\log p_{\theta}(x_t\mid x_{<t}).
\]
$\theta$ son los parámetros del modelo y $x_{<t}$ son los token previos. El objetivo solo exige aumentar la verosimilitud del siguiente token; seguir instrucciones, citar evidencia, rechazar operaciones peligrosas y planificar a largo plazo no son capacidades que este objetivo garantice directamente.
\end{importantbox}

\begin{warningbox}{El número de parámetros no es un estadístico suficiente de la capacidad}
Con la misma escala de parámetros, el número de token, la calidad de los datos, la architecture, el optimizer, la context length, el post-training y la contaminación de la evaluación modifican el resultado. Cuando más adelante la clase analiza Phi, BabyLM, DPO y MoE, lo que muestra es precisamente que, además de «más grande», existen muchos grados de libertad de diseño.
\end{warningbox}

\subsection{Emergent ability y las trampas de la medición}

La Emergent ability (capacidad emergente) se define como una capacidad que aparece en los modelos grandes y no en los pequeños. La figura de la clase muestra varias tareas cuyo desempeño sube de golpe cerca de un umbral de escala, pero el ponente también cita investigaciones que lo cuestionan: una metric discreta o no lineal puede hacer que una mejora gradual se vea como un salto. Las siguientes figuras deben leerse dos veces: la primera, para ver las curvas de las tareas; la segunda, para revisar la definición del eje vertical y el threshold de puntuación.

\lecturefigure{slide-029.jpg}{Varias tareas muestran cambios bruscos aparentes con la escala del modelo}{Diapositivas oficiales de CS25 V4, página 29}

\begin{knowledgebox}{Lectura de la figura: distinguir primero el eje horizontal, el eje vertical y el threshold}
El eje horizontal suele ser la escala del modelo o el cómputo de entrenamiento, y el vertical, la task metric. Si la metric solo otorga puntos cuando la respuesta es «completamente correcta», una pequeña mejora continua de la probabilidad subyacente puede, al cruzar el umbral, registrarse de pronto como 1. La figura puede demostrar una no linealidad en cierta evaluación, pero por sí sola no demuestra que dentro del modelo haya ocurrido una transición de fase discreta.
\end{knowledgebox}

\lecturefigure{slide-030.jpg}{Desempeño de tareas con Few-shot prompting según la escala}{Diapositivas oficiales de CS25 V4, página 30}

\begin{importantbox}{Un ejemplo sencillo de Metric-induced emergence}
Supongamos que la probabilidad de que el modelo dé la respuesta correcta aumenta de forma gradual con la escala $s$ como $p(s)$. Si la evaluación exige que $k$ pasos consecutivos sean todos correctos, la exactitud global es aproximadamente
\[
A(s)=p(s)^k.
\]
Aunque $p(s)$ sea gradual, $A(s)$ sube abruptamente en la región alta. Un análisis más sólido debe considerar a la vez la token likelihood, el partial credit, la calibration y las curvas bajo distintos umbrales.
\end{importantbox}

\lecturefigure{slide-031.jpg}{Varias explicaciones posibles de la emergent ability}{Diapositivas oficiales de CS25 V4, página 31}

\teachervoice{El ponente plantea explícitamente que la emergence puede deberse en parte a la metric no lineal que eligen los investigadores, y no a que la capacidad del modelo aparezca de verdad de forma repentina. Esto no niega el papel de la scale; exige incluir también los instrumentos de observación en el análisis causal.}

\lecturefigure{slide-032.jpg}{Beyond Scaling: la arquitectura, los datos y el proceso de entrenamiento también pueden cambiar las capacidades}{Diapositivas oficiales de CS25 V4, página 32}

\begin{knowledgebox}{Por qué beyond scaling merece un análisis aparte}
Una mejor architecture puede mejorar el sesgo inductivo; datos de mayor calidad aumentan el valor de aprendizaje de cada token; un optimizer y un curriculum más estables mejoran la eficiencia del entrenamiento; retrieval, tool y memory sacan de los parámetros parte de las capacidades. En conjunto determinan la densidad de capacidad y su accesibilidad.
\end{knowledgebox}

\lecturefigure{slide-033.jpg}{Preguntas planteadas a los estudiantes sobre scale, democratización y retrieval}{Diapositivas oficiales de CS25 V4, página 33}

\begin{warningbox}{Las preguntas de la clase no tienen una única respuesta}
«¿Basta con seguir ampliando los modelos?» y «¿Qué es mejor, la memoria paramétrica o el retrieval?» son ejes de diseño, no disyuntivas. La escala puede aportar representaciones generales, el retrieval proporciona evidencia actualizable, los modelos pequeños reducen costos y las estructuras especializadas aumentan la eficiencia. Los sistemas reales suelen combinar estas vías.
\end{warningbox}

\subsection{RLHF y DPO: de la likelihood a la preference}

El preentrenamiento hace que el modelo se ajuste a la distribución del corpus; el post-training acerca las salidas a las preferencias humanas y a las especificaciones de la tarea. El Reinforcement Learning from Human Feedback (aprendizaje por refuerzo a partir de retroalimentación humana, RLHF) suele entrenar un reward model y luego optimizar la policy con aprendizaje por refuerzo; la Direct Preference Optimization (optimización directa de preferencias, DPO) construye directamente, a partir de pares de preferencias, un objetivo de tipo clasificación.

\lecturefigure{slide-035.jpg}{Preference collection, reward model y policy optimization en RLHF}{Diapositivas oficiales de CS25 V4, página 35}

\begin{importantbox}{Forma básica de los Preference data}
Para un prompt $x$, una persona o una regla compara dos respuestas $y_w$ y $y_l$, donde $y_w$ es la preferred response. El Reward model suele aprender
\[
P(y_w\succ y_l\mid x)=\sigma\!\left(r_\phi(x,y_w)-r_\phi(x,y_l)\right).
\]
$r_\phi$ es el reward score y $\sigma$ es la logistic function. Los datos solo expresan preferencias relativas, no equivalen a la verdad objetiva; la composición de los anotadores, la rubric y la estrategia de muestreo se trasladan al comportamiento del modelo.
\end{importantbox}

\lecturefigure{slide-036.jpg}{DPO optimiza directamente el preference pair con una reference policy}{Diapositivas oficiales de CS25 V4, página 36}

\begin{importantbox}{Intuición del DPO objective}
Una forma común de escribirlo es
\[
\mathcal{L}_{\mathrm{DPO}}=-\mathbb{E}\log\sigma\!\left(\beta\left[
\log\frac{\pi_\theta(y_w\mid x)}{\pi_{\mathrm{ref}}(y_w\mid x)}-
\log\frac{\pi_\theta(y_l\mid x)}{\pi_{\mathrm{ref}}(y_l\mid x)}
\right]\right).
\]
$\pi_\theta$ es la policy por entrenar, $\pi_{\mathrm{ref}}$ es el reference model y $\beta$ controla qué tanto se permite alejarse de la reference. DPO simplifica el proceso de entrenamiento, pero el preference bias, la coverage y la over-optimization siguen presentes.
\end{importantbox}

\teachervoice{La clase no describió DPO como «RLHF ya quedó obsoleto»; subrayó que RLHF depende de retroalimentación de alta calidad, del reward y de la policy optimization, y que DPO ofrece una vía de preference-learning más directa.}

\subsection{ChatGPT, GPT-4 y Gemini: una instantánea de la clase de 2024}

Las siguientes diapositivas de modelos pertenecen al contexto del curso de 2024. Sirven para ilustrar direcciones como el post-training, la multimodalidad, el context, los benchmark y MoE, y no deben interpretarse como una clasificación de los productos más recientes de 2026. El orden de lectura debe partir de las etapas de entrenamiento y de la architecture, y después considerar los benchmark de entonces, en lugar de presuponer conclusiones a partir del nombre de la marca.

\lecturefigure{slide-037.jpg}{Panorama de ChatGPT y la serie GPT en la clase}{Diapositivas oficiales de CS25 V4, página 37}

\begin{knowledgebox}{Lectura de la figura: detrás del nombre del producto hay que separar las etapas de entrenamiento}
Un chat model suele pasar por un preentrenamiento general, un instruction tuning supervisado, preference optimization, datos de seguridad y políticas de despliegue. Que la interfaz sea parecida no significa que los datos de entrenamiento, la estructura del modelo o los límites de su comportamiento sean iguales; cuando la información pública es insuficiente, no debe deducirse todo el mecanismo interno a partir de la experiencia de uso del producto.
\end{knowledgebox}

\lecturefigure{slide-038.jpg}{Resumen de la clase sobre las capacidades, la evaluación y la multimodalidad de GPT-4}{Diapositivas oficiales de CS25 V4, página 38}

\begin{warningbox}{Un Benchmark snapshot no es una clasificación permanente}
Las puntuaciones dependen de la versión, el prompt, el conjunto de evaluación, el grado de contaminación y el judge. Los resultados presentados en clase deben conservar su fecha; no se puede usar un benchmark antiguo para afirmar que hoy cierto modelo es absolutamente superior. Lo más importante es entender las dimensiones de la evaluación: conocimiento, reasoning, multimodality, safety y cost no quedan cubiertos por una sola puntuación total.
\end{warningbox}

\lecturefigure{slide-039.jpg}{La serie Gemini y la comparación de modelos de 2024}{Diapositivas oficiales de CS25 V4, página 39}

\lecturefigure{slide-040.jpg}{La diapositiva de Gemini usa MoE para explicar los routed experts}{Diapositivas oficiales de CS25 V4, página 40}

\begin{importantbox}{Mecanismo mínimo de Mixture of Experts}
La Mixture of Experts (mezcla de expertos, MoE) usa un router para elegir unos pocos expert para cada token:
\[
p_i(x)=\operatorname{softmax}(W_rx)_i,\qquad
h'=\sum_{i\in\operatorname{TopK}(p(x))}p_i(x)E_i(x).
\]
$W_r$ son los parámetros del router, $p_i$ es el peso del expert y $E_i$ es el $i$-ésimo expert. La Sparse activation permite que, al aumentar el número total de parámetros, cada token solo se calcule con unos pocos expert; el costo es que el load balancing, la comunicación y la gestión de la capacity se vuelven más complejos.
\end{importantbox}

\begin{warningbox}{Un «Expert» no equivale a un experto humano identificable}
Un expert puede mostrar preferencia por ciertos temas o funciones, pero el entrenamiento normalmente no garantiza que se encargue solo de imágenes, matemáticas o un idioma específico. Representar MoE directamente como una división estable de oficios es solo una intuición y no sustituye el análisis empírico del routing y de la activation.
\end{warningbox}

\subsection{La situación en 2024 y las carencias por resolver}

El curso resume el estado de entonces en dos diapositivas: los LLM, la alignment, la multimodalidad y los diffusion transformer avanzaban con rapidez; los siguientes pasos incluían agent más generales, eficiencia, modelos del mundo, aprendizaje continuo y control de seguridad. En esta sección, este grupo de figuras sirve de transición entre capítulos: primero distingue las capacidades que ya se convirtieron en productos y después señala las carencias que siguen siendo objetivos de investigación, con lo que se arma una lista de problemas para las aplicaciones en otros dominios y las limitaciones de los sistemas que se tratan más adelante.

\lecturefigure{slide-041.jpg}{Where we are: síntesis de la clase de 2024 sobre el ecosistema de los LLM}{Diapositivas oficiales de CS25 V4, página 41}

\lecturefigure{slide-042.jpg}{The Future: aplicaciones más amplias, agent y eficiencia de los modelos}{Diapositivas oficiales de CS25 V4, página 42}

\begin{knowledgebox}{Lectura de la figura: la diapositiva Future mezcla tendencias técnicas y visiones de investigación}
Algunas direcciones ya cuentan con prototipos de sistemas, como la multimodal generation, el tool use y los smaller models; otras siguen siendo objetivos, como agent autónomos estables, aprendizaje continuo de bajo costo y una eficiencia de datos más cercana a la humana. Estas notas separan lo «ya demostrado» del «despliegue confiable».
\end{knowledgebox}

\lecturefigure{slide-043.jpg}{Missing ingredients: los problemas pendientes hacia sistemas más generales}{Diapositivas oficiales de CS25 V4, página 43}

\begin{importantbox}{Seis carencias tras la ampliación de capacidades}
\begin{enumerate}
  \item menor costo de entrenamiento e inferencia;
  \item representaciones más robustas del mundo real y de las relaciones causales;
  \item memoria de largo plazo actualizable y aprendizaje continuo;
  \item mecanismos internos que puedan inspeccionarse y editarse;
  \item error correction en trayectorias largas;
  \item permisos de usuario, security y value alignment.
\end{enumerate}
\end{importantbox}

\subsection{Resumen de la sección}

La Scale puede ampliar las representaciones y la cobertura de tareas, pero la emergent curve depende de la metric; la preference optimization cambia el comportamiento, pero hereda los sesgos de los datos y del objetivo; MoE, retrieval, data quality y specialized training ofrecen palancas más allá de la scale. Lo más valioso de las diapositivas de modelos no son las clasificaciones antiguas, sino la descomposición de las capacidades en etapas de entrenamiento y componentes del sistema.
